%TCIDATA{Version=5.50.0.2890}
%TCIDATA{LaTeXparent=0,0,NoticeDKE.tex}


This chapter is dedicated to the code performances. Since numerous
combinations are possible, it is not the purpose to give an exhaustive list
of results, but to select different problems that have been addressed by
existing codes in simplified magnetic configuration, a condition that allows
interesting benchmarks. Moreover, in some cases, analytical expressions are
available, and consequently accurate comparisons may be performed.

It is well known that all codes contain numerous hidden bugs, that are
difficult to track, even if a careful attention as been paid to avoid as
much as possible such problems. The only way to reduce them is to run
extensively the code for different purposes. It is the reason why both the
documentation and the code have been made as transparent as possible, so
that its robustness may be continuously improved. When some difficulties
have been encountered, they are reported in the different sections dedicated
to benchmarking. However, none of them have made the code so far unable to
determine the correct solution, providing input parameters corresponds to
assumptions used to derive the equations and their numerical counterpart. It
is well known that such kind of codes are often missused, using input
parameters well beyond the acceptable range. In some cases, warnings are
indicated by the code, but it is the duty of users to take care of this
problem before considering that the code gives a wrong answer.

The kernel of the kinetic calculations is contained by the MatLab routine
``dke\_1\_4yp.m'', whish contains both the $3-D$ zero-order Fokker-Planck
solver, and the first-order neoclassical corrections. Only this routine and
all related subroutines must be used in a chain of codes like in CRONOS
simulation package (\cite{bas03}). In order to help how to make this link,
and to benchmark the code, a test routine ``testdke\_1yp.m'' has been
written. Main global input parameters can be modified in the subroutine
``psim\_dke\_1yp.m'' that is called by ``testdke\_1yp.m'', while parameter
that are specific to a machine (Tore Supra, JET, ITER, C-MOD,\ldots ) are
gathered in the subroutine ``ptok\_dke\_1yp.m''. The machine may be virtual,
for benchmarking purposes. The input parameters for the RF wave physics
problem have be put in separate files, sinces they are quite independent of
the global input parameters. The corresponding file names for simplified RF
wave physics are called ``psim\_idealhcd\_jd.m'' or ``psim\_coldlhcd\_jd.m''
for the Lower Hybrid wave, or ``psim\_ideaeccd\_jd.m'' or
``psim\_coldeccd\_jd.m'' for the Electron Cyclotron wave. Other types of
waves are also available as well as input parameter files for more realistic
ray-tracing calculations.

Once input parameters have been set, one runs the code by executing the
MatLab script ``script\_dke\_1yp.m''. This structure allows to run several
cases in a single run, for systematic studies of the influence of one given
parameter on the final result. For example, the role played by the Ohmic
electric field may be investigated for a single value ``epsi0= 0.01 and
epsia= 0.01'', or multiple ones ``epsi0= [0.01:0.01:0.05] and epsia=
[0.01:0.01:0.05]''. In that case, all the results are saved in different
.mat binary MatLab files, in the directory that contains ``dke\_1\_4yp.m'',
under the generic name ``RESULTS\_X\_SY.mat'' where X is the machine name,
and Y the index number of the simulation. The magnetic equilibrium
calculated by HELENA is stored in an separate file named ``EQUIL\_X.mat'',
so that it may be used independently by another program for calculating
various moments of the electron distribution function like the non-thermal
bremsstrahlung emission.

All the results may be displayed using the MatLab routine
``display\_dke\_1yp.m'', with or without graphical outputs. A special
interface has been also written for an easy transfer of selected binary data
to IGOR scientific display program for final reports.

The full list of routines is given in Appendix \ref{listfiles}. Except for
arbitrary equilibrium cases that require the use of the FORTRAN F77 mexfiles
``helmex77.f'' and ``separatrix.f'' which both must be compiled\footnote{%
The code may use a simplified magnetic equilibrium with circular plasma
cross-sections, that is valid in the large aspect ration limit $\epsilon \ll
1.$ More realistic equilibria are determined numericaly. First, the external
shape of the plasma equilibrium is evaluated by specifying positions of
X-points and derivatives at this point. Then the full magnetic equilibrium
is calculated using pressure and current density profiles, in agreemnt with
the shape required at the plasma boundary. This procedure avoids numerical
instabilities that may arise from coarse estimate of the plasma shape.}, all
files are standard MatLab scripts. No external functions are used, and once
MatLab path is correctly set-up to load text files in the ``Project\_DKE''
directory, the code can be run immediatly on any computer without additional
needs.

All benchmark results discussed in the text are gathered in the file
``Benchmark.txt'', with the corresponding input parameters.
